% assets/w04-windup.svg — 와인드업은 적분기의 잘못이 아니다. 루프가 끊긴 것이다.
%
%   bash _tools/tikz2svg.sh assets/src/w04-windup.tex assets/w04-windup.svg
%
% 4주차 1-12 의 그림. 대학원 볼트 figures/src/w03-windup.tex 를 캡스톤 기호·한국어·
% 캡스톤 수치로 다시 쓴 것.
%   (a) 포화하는 동안 되먹임 루프가 끊긴다는 것을 블록선도로
%   (b) 요구가 한계 밖으로 얼마나 나갔는지를 포화 특성 위에 실척으로
%
% 기호는 1-12 본문과 같다 — tau_raw(제어기가 원한 값) · tau_sat(실제로 나간 값) ·
% K_b(되감기 이득) · 목표 y_d · 출력 y.
%
% (b) 의 수치 근거 — 1-12 유도 5번: 계단 직후 e = 1, P = K_p e = 10, tau_max = 2.5.
% 즉 P 항 하나만으로도 요구가 한계의 4 배다. (D 항까지 더한 실측 첨두 89.7 은
% 1-11-2 표에 있으나 여기서는 가로축이 너무 길어져 P 항만 표시한다.)
%
% 배치 주의 — 설명 글은 그림 위에 얹지 않는다. (a) 는 아래 한 줄, (b) 는 오른쪽 한 칸에
% 모아 둔다. 처음에는 라벨을 그림 안에 넣었다가 배선·축과 전부 겹쳤다.
%
% TeX 제어어 이름에는 숫자를 못 쓴다. \def\Yz 처럼 글자만 쓴다.
%
% 폭 약 148 mm. 본문 폭 186 mm 안.
\documentclass[border=6pt]{standalone}
% capstone-style.tex — 이 볼트의 TikZ 그림이 공유하는 스타일
%
% 쓰는 법:  % capstone-style.tex — 이 볼트의 TikZ 그림이 공유하는 스타일
%
% 쓰는 법:  % capstone-style.tex — 이 볼트의 TikZ 그림이 공유하는 스타일
%
% 쓰는 법:  \input{capstone-style}   (assets/src/ 안에서)
% 빌드:     bash _tools/tikz2svg.sh assets/src/w02-node-topic.tex assets/w02-node-topic.svg
%
% 왜 TikZ 인가
%   손으로 SVG 좌표를 쓰면 화살촉·선 굵기·글자 크기가 그림마다 미묘하게 달라진다.
%   그것이 "자료가 어설퍼 보인다"의 정체다. 스타일을 한 번만 정하고 상대좌표로
%   배치하면 좌표를 고쳐도 그림이 무너지지 않는다.
%   수식은 본문과 같은 TeX 글꼴로 나온다.
%
% 한글
%   ko.TeX(DVI 경로)를 쓴다. XeLaTeX 은 TikZ 그래픽을 PDF special 로 내보내는데
%   dvisvgm 이 그것을 받으려면 Ghostscript < 10.01 이 필요하다. 이 PC 에는
%   10.07 이 깔려 있어 그 경로가 막힌다. latex -> DVI -> dvisvgm 이 유일하게
%   동작하는 조합이다.

\usepackage[hangul]{kotex}
\usepackage{tikz}
\usepackage{amsmath}
\usepackage{xcolor}
\usetikzlibrary{arrows.meta, positioning, calc, fit, backgrounds, decorations.pathmorphing}

% ---- 색 --------------------------------------------------------------------
% 강조는 하나만. 나머지는 검정. 흑백 인쇄에서도 읽혀야 한다.
\definecolor{ink}{HTML}{1A1A1A}
\definecolor{soft}{HTML}{666666}
\definecolor{accent}{HTML}{7C3AED}
\definecolor{warn}{HTML}{B91C1C}

% 계층 색 — Simulink 모델의 블록 색과 같은 규칙을 쓴다
\definecolor{guid}{HTML}{1D4ED8}    % 유도
\definecolor{ctrl}{HTML}{EA580C}    % 제어
\definecolor{thr}{HTML}{D97706}     % 추진기
\definecolor{plant}{HTML}{16A34A}   % 운동모델
\definecolor{nav}{HTML}{7C3AED}     % 항법 · 신호처리
\definecolor{pathc}{HTML}{0D9488}   % 계획 경로 · 항적
\definecolor{errc}{HTML}{C2410C}    % 오차
\definecolor{flow}{HTML}{0369A1}    % 조류 · 바람

% ---- 선과 화살촉 -----------------------------------------------------------
\tikzset{
  sig/.style   = {ink, line width=0.5pt, -{Stealth[length=4pt, width=3pt]}},
  wire/.style  = {ink, line width=0.5pt},
  fb/.style    = {ink, line width=0.5pt, densely dashed,
                  -{Stealth[length=4pt, width=3pt]}},
  asig/.style  = {accent, line width=0.5pt, densely dashed,
                  -{Stealth[length=4pt, width=3pt]}},
  % 블록 — 흰 바탕, 직각, 검은 테두리. 색은 테두리에만 준다
  blk/.style   = {draw=ink, line width=0.55pt, fill=white,
                  minimum width=18mm, minimum height=8.5mm, inner sep=2pt, font=\small},
  gblk/.style  = {blk, draw=guid},
  cblk/.style  = {blk, draw=ctrl},
  tblk/.style  = {blk, draw=thr},
  pblk/.style  = {blk, draw=plant},
  nblk/.style  = {blk, draw=nav},
  % 합산점 · 분기점
  sum/.style   = {draw=ink, line width=0.55pt, fill=white, circle,
                  inner sep=0pt, minimum size=4.6mm, font=\scriptsize},
  dot/.style   = {ink, fill=ink, circle, inner sep=0pt, minimum size=1.5mm},
  % 글자
  sname/.style = {font=\small, inner sep=1.5pt},
  note/.style  = {font=\scriptsize, soft, inner sep=1.5pt},
  wnote/.style = {font=\scriptsize, warn, inner sep=1.5pt},
  anote/.style = {font=\scriptsize, accent, inner sep=1.5pt},
  axis/.style  = {ink, line width=0.5pt},
  tick/.style  = {font=\scriptsize, soft},
  % 묶음 상자 — 서브시스템 표시
  grp/.style   = {draw=soft, line width=0.4pt, densely dotted, rounded corners=1mm,
                  inner sep=3mm},
}
   (assets/src/ 안에서)
% 빌드:     bash _tools/tikz2svg.sh assets/src/w02-node-topic.tex assets/w02-node-topic.svg
%
% 왜 TikZ 인가
%   손으로 SVG 좌표를 쓰면 화살촉·선 굵기·글자 크기가 그림마다 미묘하게 달라진다.
%   그것이 "자료가 어설퍼 보인다"의 정체다. 스타일을 한 번만 정하고 상대좌표로
%   배치하면 좌표를 고쳐도 그림이 무너지지 않는다.
%   수식은 본문과 같은 TeX 글꼴로 나온다.
%
% 한글
%   ko.TeX(DVI 경로)를 쓴다. XeLaTeX 은 TikZ 그래픽을 PDF special 로 내보내는데
%   dvisvgm 이 그것을 받으려면 Ghostscript < 10.01 이 필요하다. 이 PC 에는
%   10.07 이 깔려 있어 그 경로가 막힌다. latex -> DVI -> dvisvgm 이 유일하게
%   동작하는 조합이다.

\usepackage[hangul]{kotex}
\usepackage{tikz}
\usepackage{amsmath}
\usepackage{xcolor}
\usetikzlibrary{arrows.meta, positioning, calc, fit, backgrounds, decorations.pathmorphing}

% ---- 색 --------------------------------------------------------------------
% 강조는 하나만. 나머지는 검정. 흑백 인쇄에서도 읽혀야 한다.
\definecolor{ink}{HTML}{1A1A1A}
\definecolor{soft}{HTML}{666666}
\definecolor{accent}{HTML}{7C3AED}
\definecolor{warn}{HTML}{B91C1C}

% 계층 색 — Simulink 모델의 블록 색과 같은 규칙을 쓴다
\definecolor{guid}{HTML}{1D4ED8}    % 유도
\definecolor{ctrl}{HTML}{EA580C}    % 제어
\definecolor{thr}{HTML}{D97706}     % 추진기
\definecolor{plant}{HTML}{16A34A}   % 운동모델
\definecolor{nav}{HTML}{7C3AED}     % 항법 · 신호처리
\definecolor{pathc}{HTML}{0D9488}   % 계획 경로 · 항적
\definecolor{errc}{HTML}{C2410C}    % 오차
\definecolor{flow}{HTML}{0369A1}    % 조류 · 바람

% ---- 선과 화살촉 -----------------------------------------------------------
\tikzset{
  sig/.style   = {ink, line width=0.5pt, -{Stealth[length=4pt, width=3pt]}},
  wire/.style  = {ink, line width=0.5pt},
  fb/.style    = {ink, line width=0.5pt, densely dashed,
                  -{Stealth[length=4pt, width=3pt]}},
  asig/.style  = {accent, line width=0.5pt, densely dashed,
                  -{Stealth[length=4pt, width=3pt]}},
  % 블록 — 흰 바탕, 직각, 검은 테두리. 색은 테두리에만 준다
  blk/.style   = {draw=ink, line width=0.55pt, fill=white,
                  minimum width=18mm, minimum height=8.5mm, inner sep=2pt, font=\small},
  gblk/.style  = {blk, draw=guid},
  cblk/.style  = {blk, draw=ctrl},
  tblk/.style  = {blk, draw=thr},
  pblk/.style  = {blk, draw=plant},
  nblk/.style  = {blk, draw=nav},
  % 합산점 · 분기점
  sum/.style   = {draw=ink, line width=0.55pt, fill=white, circle,
                  inner sep=0pt, minimum size=4.6mm, font=\scriptsize},
  dot/.style   = {ink, fill=ink, circle, inner sep=0pt, minimum size=1.5mm},
  % 글자
  sname/.style = {font=\small, inner sep=1.5pt},
  note/.style  = {font=\scriptsize, soft, inner sep=1.5pt},
  wnote/.style = {font=\scriptsize, warn, inner sep=1.5pt},
  anote/.style = {font=\scriptsize, accent, inner sep=1.5pt},
  axis/.style  = {ink, line width=0.5pt},
  tick/.style  = {font=\scriptsize, soft},
  % 묶음 상자 — 서브시스템 표시
  grp/.style   = {draw=soft, line width=0.4pt, densely dotted, rounded corners=1mm,
                  inner sep=3mm},
}
   (assets/src/ 안에서)
% 빌드:     bash _tools/tikz2svg.sh assets/src/w02-node-topic.tex assets/w02-node-topic.svg
%
% 왜 TikZ 인가
%   손으로 SVG 좌표를 쓰면 화살촉·선 굵기·글자 크기가 그림마다 미묘하게 달라진다.
%   그것이 "자료가 어설퍼 보인다"의 정체다. 스타일을 한 번만 정하고 상대좌표로
%   배치하면 좌표를 고쳐도 그림이 무너지지 않는다.
%   수식은 본문과 같은 TeX 글꼴로 나온다.
%
% 한글
%   ko.TeX(DVI 경로)를 쓴다. XeLaTeX 은 TikZ 그래픽을 PDF special 로 내보내는데
%   dvisvgm 이 그것을 받으려면 Ghostscript < 10.01 이 필요하다. 이 PC 에는
%   10.07 이 깔려 있어 그 경로가 막힌다. latex -> DVI -> dvisvgm 이 유일하게
%   동작하는 조합이다.

\usepackage[hangul]{kotex}
\usepackage{tikz}
\usepackage{amsmath}
\usepackage{xcolor}
\usetikzlibrary{arrows.meta, positioning, calc, fit, backgrounds, decorations.pathmorphing}

% ---- 색 --------------------------------------------------------------------
% 강조는 하나만. 나머지는 검정. 흑백 인쇄에서도 읽혀야 한다.
\definecolor{ink}{HTML}{1A1A1A}
\definecolor{soft}{HTML}{666666}
\definecolor{accent}{HTML}{7C3AED}
\definecolor{warn}{HTML}{B91C1C}

% 계층 색 — Simulink 모델의 블록 색과 같은 규칙을 쓴다
\definecolor{guid}{HTML}{1D4ED8}    % 유도
\definecolor{ctrl}{HTML}{EA580C}    % 제어
\definecolor{thr}{HTML}{D97706}     % 추진기
\definecolor{plant}{HTML}{16A34A}   % 운동모델
\definecolor{nav}{HTML}{7C3AED}     % 항법 · 신호처리
\definecolor{pathc}{HTML}{0D9488}   % 계획 경로 · 항적
\definecolor{errc}{HTML}{C2410C}    % 오차
\definecolor{flow}{HTML}{0369A1}    % 조류 · 바람

% ---- 선과 화살촉 -----------------------------------------------------------
\tikzset{
  sig/.style   = {ink, line width=0.5pt, -{Stealth[length=4pt, width=3pt]}},
  wire/.style  = {ink, line width=0.5pt},
  fb/.style    = {ink, line width=0.5pt, densely dashed,
                  -{Stealth[length=4pt, width=3pt]}},
  asig/.style  = {accent, line width=0.5pt, densely dashed,
                  -{Stealth[length=4pt, width=3pt]}},
  % 블록 — 흰 바탕, 직각, 검은 테두리. 색은 테두리에만 준다
  blk/.style   = {draw=ink, line width=0.55pt, fill=white,
                  minimum width=18mm, minimum height=8.5mm, inner sep=2pt, font=\small},
  gblk/.style  = {blk, draw=guid},
  cblk/.style  = {blk, draw=ctrl},
  tblk/.style  = {blk, draw=thr},
  pblk/.style  = {blk, draw=plant},
  nblk/.style  = {blk, draw=nav},
  % 합산점 · 분기점
  sum/.style   = {draw=ink, line width=0.55pt, fill=white, circle,
                  inner sep=0pt, minimum size=4.6mm, font=\scriptsize},
  dot/.style   = {ink, fill=ink, circle, inner sep=0pt, minimum size=1.5mm},
  % 글자
  sname/.style = {font=\small, inner sep=1.5pt},
  note/.style  = {font=\scriptsize, soft, inner sep=1.5pt},
  wnote/.style = {font=\scriptsize, warn, inner sep=1.5pt},
  anote/.style = {font=\scriptsize, accent, inner sep=1.5pt},
  axis/.style  = {ink, line width=0.5pt},
  tick/.style  = {font=\scriptsize, soft},
  % 묶음 상자 — 서브시스템 표시
  grp/.style   = {draw=soft, line width=0.4pt, densely dotted, rounded corners=1mm,
                  inner sep=3mm},
}

\tikzset{
  asum/.style  = {draw=accent, line width=0.55pt, fill=white, circle,
                  inner sep=0pt, minimum size=4.6mm, font=\scriptsize},
  awire/.style = {accent, line width=0.5pt},
}

\begin{document}
\begin{tikzpicture}[x=1mm, y=1mm]

% ===================== (a) 루프가 끊기는 자리 ===============================
\begin{scope}
  \node[font=\small\bfseries, anchor=west] at (-14, 20) {(a) \ \ 포화하는 동안 루프가 끊긴다};

  \node[sum]                                (sa)  at (0,0) {};
  \node[cblk, right=9mm of sa, yshift=9mm, minimum width=14mm]  (kp)  {$K_p$};
  \node[cblk, right=9mm of sa, yshift=-9mm, minimum width=14mm] (ki)  {$K_i/s$};
  \node[sum,  right=42mm of sa]             (sb)  {};
  \node[tblk, right=15mm of sb, minimum width=14mm] (sat) {};
  \node[pblk, right=15mm of sat, minimum width=16mm] (pl)  {플랜트};

  % 포화 글리프 — 상자 안에 글자를 또 넣지 않는다
  \draw[wire] ([shift={(-5.5mm,-2.4mm)}]sat.center) -- ([shift={(-2mm,-2.4mm)}]sat.center)
              -- ([shift={( 2mm, 2.4mm)}]sat.center) -- ([shift={( 5.5mm, 2.4mm)}]sat.center);
  \node[note, below=0.8mm of sat] {포화 \textsf{Sat}};

  % 입력과 분기
  \draw[sig] (-14,0) -- node[sname, above, pos=0.22] {$y_d$} (sa);
  \coordinate (br) at ($(sa)+(4.5mm,0)$);
  \draw[wire] (sa) -- (br);
  \node[dot] at (br) {};
  \draw[wire] (br) |- ($(kp.west)+(-3mm,0)$);
  \draw[sig]  ($(kp.west)+(-3mm,0)$) -- (kp);
  \draw[wire] (br) |- ($(ki.west)+(-3mm,0)$);
  \draw[sig]  ($(ki.west)+(-3mm,0)$) -- (ki);
  \node[sname, above right=0mm and -0.5mm of br] {$e$};

  % 합류
  \coordinate (jn) at ($(sb)+(-5.5mm,0)$);
  \draw[wire] (kp) -| (jn);  \draw[wire] (ki) -| (jn);
  \draw[sig]  (jn) -- (sb);

  % 앞으로 가는 길
  \draw[sig] (sb)  -- node[sname, above] {$\tau_{\text{raw}}$} (sat);
  \draw[sig] (sat) -- node[sname, above, pos=0.26] {$\tau_{\text{sat}}$} (pl);
  \coordinate (out) at ($(pl.east)+(10mm,0)$);
  \draw[wire] (pl) -- (out);
  \draw[sig]  (out) -- ++(8mm,0) node[sname, above, pos=0.5] {$y$};

  % 부호
  \node[font=\scriptsize] at ($(sa)+(-2.8mm, 3.0mm)$) {$+$};
  \node[font=\scriptsize] at ($(sa)+(-3.6mm,-3.2mm)$) {$-$};
  \node[font=\scriptsize] at ($(sb)+(-3.0mm, 3.0mm)$) {$+$};
  \node[font=\scriptsize] at ($(sb)+( 3.0mm, 3.0mm)$) {$+$};

  % 되먹임
  \node[dot] at (out) {};
  \draw[wire] (out) |- (0,-26);
  \draw[sig]  (0,-26) -- (sa);

  % 루프가 끊기는 자리 — 표시는 선 위에, 글은 위쪽 빈 곳에
  \coordinate (cut) at ($(sat.east)!0.80!(pl.west)$);
  \draw[warn, line width=0.8pt] ($(cut)+(-1.6mm,-1.6mm)$) -- ($(cut)+(1.6mm,1.6mm)$);
  \draw[warn, line width=0.8pt] ($(cut)+(-1.6mm, 1.6mm)$) -- ($(cut)+(1.6mm,-1.6mm)$);
  \node[wnote, anchor=south, align=center, text width=46mm] at ($(cut)+(0,6.4mm)$)
       {포화 중에는 여기가 끊긴다 —\\ 게인을 더 올려도 나가는 값이 늘지 않는다};

  % 되감기 — SumAW 와 Kb_gain 이 이 자리다
  \coordinate (tc) at ($(sb)!0.42!(sat.west)$);
  \coordinate (ts) at ($(sat.east)!0.26!(pl.west)$);
  \node[dot] at (tc) {};  \node[dot] at (ts) {};
  \node[asum] (aw) at ($(tc)!0.5!(ts) + (0,-15mm)$) {};
  \draw[awire] (tc) |- ($(aw.west)+(-4.5mm,0)$);
  \draw[asig]  ($(aw.west)+(-4.5mm,0)$) -- (aw);
  \draw[awire] (ts) |- ($(aw.east)+(4.5mm,0)$);
  \draw[asig]  ($(aw.east)+(4.5mm,0)$) -- (aw);
  \node[font=\scriptsize, accent] at ($(aw)+(-3.2mm,3.2mm)$) {$-$};
  \node[font=\scriptsize, accent] at ($(aw)+( 3.2mm,3.2mm)$) {$+$};
  \node[anote, anchor=south] at ($(aw)+(0,5.4mm)$) {$\tau_{\text{sat}} - \tau_{\text{raw}}$};
  \draw[awire] (aw) -- ++(0,-5mm) -| ($(ki.south)+(0,-5mm)$);
  \draw[asig]  ($(ki.south)+(0,-5mm)$) -- (ki.south);

  \node[anote, anchor=north west, align=left, text width=150mm] at (-14,-30)
       {보라색이 되감기 —\ \textsf{PID\_byhand} 안의 \textsf{SumAW} 가 이 뺄셈이고
        \textsf{Kb\_gain} 이 그 차이에 $K_b$ 를 곱해 적분기 입력으로 되돌린다.
        포화가 풀리면 두 값이 같아져 이 갈래가 저절로 0 이 된다};
\end{scope}

% ===================== (b) 한계 밖으로 얼마나 나갔는가 ======================
\begin{scope}[shift={(6,-86)}]
  \node[font=\small\bfseries, anchor=west] at (-18, 32)
       {(b) \ \ 요구가 한계 밖으로 얼마나 나갔는가 \ \ (1-12 의 조건, 계단 직후)};

  % 자료 -> mm.  x: tau_raw 1 단위 = 5 mm ,  y: tau_sat 1 단위 = 3.25 mm
  \def\Xs{5.0}
  \def\Ys{3.25}
  \def\Xz{10.0}          % tau_raw = 0 의 가로 위치
  \def\Yz{13.0}          % tau_sat = 0 의 세로 위치

  \draw[axis] (-1,\Yz) -- (70,\Yz);
  \draw[axis] (\Xz,-1) -- (\Xz,26);

  % 포화 특성 — -2.5 에서 꺾이고 2.5 에서 다시 꺾인다
  \draw[ink, line width=0.9pt]
        (-1, {\Yz - 2.5*\Ys}) -- ({\Xz - 2.5*\Xs}, {\Yz - 2.5*\Ys})
     -- ({\Xz + 2.5*\Xs}, {\Yz + 2.5*\Ys}) -- (68, {\Yz + 2.5*\Ys});
  \draw[soft, line width=0.35pt, densely dotted]
        (\Xz, {\Yz + 2.5*\Ys}) -- (68, {\Yz + 2.5*\Ys});

  % 요구값 10 — 1-12 유도 5번의 P 항
  \draw[accent, line width=0.5pt, densely dashed]
        ({\Xz + 10*\Xs}, \Yz) -- ({\Xz + 10*\Xs}, {\Yz + 2.5*\Ys});
  \fill[accent] ({\Xz + 10*\Xs}, {\Yz + 2.5*\Ys}) circle (0.8mm);
  % 버려지는 양을 가로 화살표로
  \draw[accent, line width=0.5pt, {Stealth[length=3.5pt,width=2.6pt]}-{Stealth[length=3.5pt,width=2.6pt]}]
        ({\Xz + 2.5*\Xs}, {\Yz + 2.5*\Ys + 3.0}) -- ({\Xz + 10*\Xs}, {\Yz + 2.5*\Ys + 3.0});
  \node[anote, anchor=south] at ({\Xz + 6.25*\Xs}, {\Yz + 2.5*\Ys + 3.4}) {버려지는 $7.5$};

  % 가로 눈금
  \foreach \v in {-2, 0, 2.5} {
    \draw[axis] ({\Xz + \v*\Xs}, \Yz) -- ({\Xz + \v*\Xs}, \Yz-1.2);
    \node[tick, below=0.3mm] at ({\Xz + \v*\Xs}, \Yz-1.2) {\v};
  }
  \draw[accent, line width=0.5pt] ({\Xz + 10*\Xs}, \Yz) -- ({\Xz + 10*\Xs}, \Yz-1.2);
  \node[anote, below=0.3mm] at ({\Xz + 10*\Xs}, \Yz-1.2) {$10$};
  \node[font=\small, anchor=north] at (35, \Yz-7) {$\tau_{\text{raw}}$ \ \ 제어기가 원한 값};
  % 세로 눈금 — 한계 두 개만
  \foreach \v in {2.5, -2.5} {
    \draw[axis] (\Xz, {\Yz + \v*\Ys}) -- (\Xz-1.3, {\Yz + \v*\Ys});
    \node[tick, anchor=east] at (\Xz-1.8, {\Yz + \v*\Ys}) {\v};
  }
  \node[font=\small, anchor=south, rotate=90] at (-6, \Yz) {$\tau_{\text{sat}}$ \ \ 실제로 나간 값};

  % 오른쪽 설명 칸 — 그림 위에 글을 얹지 않는다
  \node[note, anchor=north west, align=left, text width=54mm] at (74, 29)
       {가운데 기울기 $1$ 인 구간에서만 $\tau_{\text{sat}} = \tau_{\text{raw}}$ 다.
        평평한 두 구간이 포화다.\\[2mm]
        계단 직후 $e = 1$ 이므로 $P$ 항 하나만으로 $\tau_{\text{raw}} = K_p e = 10$ —
        한계 $2.5$ 의 네 배다 (1-12 유도 5번).\\[2mm]
        \textbf{이 그림에는 적분기가 없다.} 그런데도 요구와 실제가 다르다.
        와인드업은 적분기가 만드는 것이 아니라, 이 간격이 벌어져 있는 동안
        적분기가 그것을 모른 채 계속 쌓아서 생긴다};

\end{scope}

\end{tikzpicture}
\end{document}
